\section{La perspectiva del CEO: el periodo de incertidumbre, la energía mental y el All In}

Otro aspecto importante de este episodio es cómo un emprendedor gestiona el periodo de exploración. Muyan Zhiyu (沐言智语) obtuvo financiamiento sin contratiempos, pero Zhang Yueguang (张月光) repitió varias veces que el dinero no es lo más importante: lo es más la energía mental del equipo y la suya propia. Una exploración sin resultados desgasta al equipo; si el fundador cambia de dirección una y otra vez sin gestionar las expectativas, el equipo termina por derrumbarse.

\subsection{Cómo comunicarse en el periodo de incertidumbre}

Describe los principios que sigue en el periodo de incertidumbre: atenerse a los hechos y no presentar como dirección segura una idea en la que no tiene confianza; decirle al equipo que todavía no se ha hecho all in; permitir que otros propongan mejores ideas; y probar con errores de bajo costo para cuidar la energía mental del equipo. Esta forma de actuar coincide con su método de exploración de producto: cuando hay incertidumbre, se controlan las variables; solo cuando hay certeza se concentran los recursos.

\begin{knowledgebox}{Gestión organizacional en el periodo de exploración}
En un emprendimiento de aplicaciones de IA es fácil caer en un estado en el que «cada dirección parece una oportunidad». En el periodo de exploración hay que proteger el efectivo y, más aún, la energía mental del equipo; antes de hacer all in, hay que decirle claramente al equipo que se trata solo de una prueba, no de la batalla definitiva.
\end{knowledgebox}

\subsection{Por qué con Dokie sí se puede hacer All In}

Zhang Yueguang cuenta que, cuando pensó en lo que hay detrás de Dokie ---crear un amigo de IA que rebase los límites de sus propias capacidades---, aceptaba incluso perder el dinero de la empresa en ese proyecto. Es un criterio exigente: no se trata de si un producto aislado de PPT puede volverse un éxito, sino de si la gran tesis que hay detrás merece que la empresa apueste por ella.

\begin{importantbox}{La decisión de hacer All In}
El all in no debe surgir de la ansiedad, de la presión del financiamiento ni de perseguir la moda, sino de la convicción del fundador sobre una gran tesis: aun si fracasa, sabe por qué valía la pena perder en ella.
\end{importantbox}

\subsection{Resumen de la sección}

Una empresa de aplicaciones de IA no solo tiene que encontrar un producto, también tiene que gestionar el periodo de exploración. El all in de Zhang Yueguang con Dokie es, en realidad, una apuesta por «un amigo de IA que rebase los límites de sus capacidades», y no por la herramienta puntual de PPT.
